This appendix records the implementation choices needed to interpret the main-text comparisons.

\paragraph{Domains.}
The two main exact-checker domains are arithmetic and linear equations.
Both domains support prefix-level oracle evaluation through rollout under a fixed base reasoner and exact verification.
This is what makes action-level utility, determinacy, and crossing measurable at intermediate prefixes rather than only at final answers.

\paragraph{Utility, abstention, and ambiguity.}
For each non-abstain action, the oracle estimates expected utility from a small Monte Carlo rollout set under a fixed base reasoner.
The rollout utility is
\[
\mathbb{1}\{\text{success}\} - \lambda_{\text{tok}} \cdot \text{tokens} - \gamma_{\text{wrong}} \cdot \mathbb{1}\{\text{unsafe wrong}\},
\]
while \abstain is assigned utility \(0\).
Across the benchmark datasets used in the paper, each non-abstain action is estimated from \(4\)--\(8\) rollouts.
A prefix is marked ambiguous if the top-two action estimates are not cleanly separated, either because the mean-utility gap is below \(0.05\) or because the corresponding \(95\%\) confidence intervals overlap.
Main-text controller comparisons exclude ambiguous prefixes by default.

\paragraph{Prefix sampling and bounded revision.}
Decision points are sampled along base traces produced by the fixed reasoner, and recoverable perturbations are injected when needed to cover revise-relevant states.
The \revise action is intentionally local: it rolls back the most recent faulty step and resumes rollout from the repaired prefix.
This bounded operator is part of the benchmark design.
It makes revision measurable at the prefix level without turning the oracle into an open-ended planner over arbitrary editing strategies.

\paragraph{Frozen main-text controller set.}
The main paper freezes the controller comparison to five rows:
\orderedscalar, \learnedoned, \directpolicy, exact-state \triver, and one selected deployable predicted-state \triver controller per domain.
This freeze is intentional.
The paper is about benchmark design and controller-class comparison, not about exhaustively searching over every local architecture variant.
Broader value-head and routing sweeps are therefore treated as appendix-only evidence.

\paragraph{Selected deployable predicted-state controller.}
The deployable predicted-state row is domain-specific.
For linear equations, the selected controller uses the predicted state together with an exact-plus-out-of-fold value-head training regime.
For arithmetic, the selected controller uses the strongest available exact-trained deployable predicted-state controller.
In the main text we report this row uniformly as predicted-state \triver, because the conceptual question is whether the structured controller survives deployment, not which local variant happened to be the strongest within a domain.
The later appendix rescue probes intentionally do not overwrite this row.
They are used to diagnose whether the remaining gap is due to label noise, target mismatch, or state identification, not to reopen the main-text leaderboard.

\paragraph{Repeatability protocol.}
The repeatability table in the main text uses the \(100\)-sample scale and reports means and standard deviations over a minimal set of within-domain reruns.
Its purpose is not to certify a universal winner.
Its purpose is to constrain wording: if a ranking does not survive this repeatability check, we do not write it as a stable paper claim.
